% Neutral Article demonstration, included by nature-cardiovascular-research.tex.
% The manuscript example, author labels, diagram and references are synthetic.
\begingroup
\newcommand{\NCRArticleType}{\sffamily\bfseries\fontsize{9bp}{11bp}\selectfont 研究论文\enspace Article}
\newcommand{\NCRArticleCite}[1]{\textsuperscript{\hyperlink{ncv-article:ref:#1}{#1}}}
\newcommand{\NCRArticleFigure}{\hyperlink{ncv-article:fig:1}{图~1}}
\fancyhead[L]{\NCRArticleType}
\fancyhead[R]{\LetterIdentifier}
\fancyfoot[L]{\sffamily\fontsize{7.2bp}{8.6bp}\selectfont\textcolor{LetterBlue}{Nature Cardiovascular Research}}

\thispagestyle{letteropening}
\vspace*{-25bp}
{\color{LetterPink}\hrule height 3.5bp}\vspace{5bp}
\noindent{\sffamily\bfseries\fontsize{17bp}{21bp}\selectfont Nature Cardiovascular Research\par}
\vspace{37bp}
\noindent{\NCRArticleType}\hfill{\LetterIdentifier}\par
\vspace{3bp}\hrule height .35bp\vspace{5bp}
\noindent{\bfseries\fontsize{26bp}{32bp}\selectfont
以多通道观测记录为线索\par
\noindent 识别馆藏展示条件的空间差异\par
\noindent 并复核跨时段对应关系\par}
\vspace{23bp}\hrule height .35bp\vspace{6bp}

\noindent\begin{minipage}[t]{.317\textwidth}
  \vspace{0pt}\sffamily\fontsize{8bp}{11.3bp}\selectfont
  收稿日期\enspace 按原文填写\par\vspace{4bp}\hrule height .3bp\vspace{4bp}
  接受日期\enspace 按原文填写\par\vspace{4bp}\hrule height .3bp\vspace{4bp}
  在线发表\enspace 原稿尚未填写\par\vspace{4bp}\hrule height .3bp\vspace{4bp}
  非官方中文样板
\end{minipage}\hfill
\begin{minipage}[t]{.655\textwidth}
  \vspace{0pt}\setlength{\parindent}{0pt}
  {\sffamily\bfseries\fontsize{9bp}{12.1bp}\selectfont
  示例作者甲\textsuperscript{1}、示例作者乙\textsuperscript{2}、示例作者丙\textsuperscript{1,2}、
  示例作者丁\textsuperscript{3}、示例作者戊\textsuperscript{1}、示例作者己\textsuperscript{2}、
  示例作者庚\textsuperscript{3}、示例作者辛\textsuperscript{1,3}、示例作者壬\textsuperscript{2}\par}
  \vspace{6bp}\hrule height .35bp\vspace{7bp}
  {\fontsize{11bp}{15.2bp}\selectfont
  展室记录通常同时包含位置、时段和观测条件。将三个维度保存在同一记录中，便于复查陈列调整与不同读数之间的关系。这个中性示例采用可重复的记录流程，将展示条件与观测表逐项对应，并为每个时段保留起止信息。示意图说明多个记录入口如何汇入同一份核查表。图中的位置、标记和连线均用于说明排版，无实际馆藏数据。研究论文译稿可在此处保留连续摘要，在同一页接排双栏正文，并在首页脚注提示作者单位的文末位置。\par}
\end{minipage}
\par\vspace{10bp}

\begin{multicols}{2}
\noindent 展室记录可按房间划分，也可按采集时段排列。若观测人员在不同时段修改了记录方式，两次记录即使使用同一地点名称，也可能对应不同的观察条件。复查时应先找出每条记录的时段和来源，再确定它能否与其他条目成组使用\NCRArticleCite{1}。

同一展室内的记录需要保留足够的位置信息。例如，可用简图说明窗侧、过道和展示区之间的相对位置，而把具体采集时间留在记录表。简图中移动标记的位置可能改变读者对记录范围的理解，因此图注应说明节点、箭头和门窗标记的实际含义。

时段信息来自原始观测表，每个时段应能够回查到相应版本。若两个观察时段有重叠，可在记录中保留重复部分，标示两次采集的范围。只知道先后顺序时，可使用明确的顺序词，并说明准确时间未载明。补录的内容可另用变更记录注明。

本示例将来源表、整理表和核查表作为连续的阅读环节。先在来源表辨别记录，再经整理表核实版本，最后在核查表确认字段之间的关系。\NCRArticleFigure{}展示这些步骤；它保留完整图注，图和紧随其后的正文可在同页阅读。

跨时段对应关系还可以用于核对记录缺项。若某一时段没有观察结果，核查表应区分未采集、无法定位和待核实等状态，并保留作出判断的依据。不能凭另一时段已有读数推补空白。后续取得材料时，可沿用现有记录编号补充来源说明\NCRArticleCite{2}。

不同栏目采用的术语也需要对应。例如，位置名称可能来自馆区地图，记录时段可能来自值班表。将两套材料合并之前，应分别核对资料范围，并记录每个名称适用的时段。若次年更换名称，可在整理表中并列保留两个版本。

\subsection*{记录边界}
记录中的一条线可表示位置的边界，另一条线可能表示两项材料之间的索引关系。它们的含义在图例中应分别说明。后续增补注释时，应让文字注释关联原有的符号，保留能定位原图位置的信息。

各观察区可用同一套编号登记，具体页次则随文件版本记录。若某条记录在多个文件中出现，可注明哪一份最早以及后来是否作过修订。合并视图可用稳定的内部编号，将同一位置的记录排在一起。

\subsection*{记录单位}
核查表需要注明每一行代表的位置或时段。一个位置若设几处观察点，可分别列出各观察点；一条记录若覆盖整个展室，应写明其代表范围。通过这些字段可从核查表重新定位相应来源。

如果资料中列出位置、时间和采集条件，译稿应保留它们的对应关系。条件说明可跟随记录一同阅读；条目过长时，可以利用跨行分组和小标题标记边界。图文长短不一时，优先让图注和可续排的双栏文字共同利用页面。

\subsection*{来源记录的复核}
可依次检查记录编号、对应位置和适用时段，并在有疑问的条目旁标明待核实字段。后续发现排版遗漏时，可根据原编号回到相关图或表，找到原件中与译文对应的项目。

资料中对记录范围作了限定时，整理表应注明该限定适用于哪些位置或哪些时段。没有发生变化的字段可继续使用原来的编号，新增字段则注明首次纳入整理表的版本。
\end{multicols}

\vfill\hrule height .35bp\vspace{5bp}
\noindent{\sffamily\fontsize{7.65bp}{10bp}\selectfont
完整作者单位见\hyperlink{ncv-article:affiliations}{文末}。\enspace 通信信息按所给原稿填写。\par}

\clearpage
\noindent\begin{minipage}{\textwidth}
\centering\hypertarget{ncv-article:fig:1}{}
\begin{tikzpicture}[font=\sffamily\small,>=Stealth,
  box/.style={draw=LetterBlue,fill=LetterBlue!7,rounded corners=2bp,align=center,
    text width=108bp,minimum height=29bp,inner sep=4bp}]
  \node[box] (location) at (0,1.55) {位置索引\strut\\展室与观察点};
  \node[box] (time) at (0,-1.55) {时段索引\strut\\起止时间与记录来源};
  \node[box] (prepare) at (5.87,0) {整理表\strut\\统一字段和版本标记};
  \node[box,draw=LetterPink!85!black,fill=LetterPink!6] (review) at (11.74,0) {核查表\strut\\确认可对应的记录};
  \draw[->,thick] (location.east)--(prepare.west);
  \draw[->,thick] (time.east)--(prepare.west);
  \draw[->,thick] (prepare.east)--(review.west);
\end{tikzpicture}\par\vspace{4bp}
\raggedright\sffamily\fontsize{7.8bp}{10.25bp}\selectfont
\textbf{图~1｜从来源到核查的记录流程。}位置索引与时段索引分别提供原始字段。两组字段汇入整理表，版本复核完成后写入核查表。
箭头表示阅读与核对的顺序。图中的标签、位置和连线为原创版式占位。\par
\end{minipage}\par\medskip

\begin{multicols}{2}
\section*{整理和复查}
核查从逐项对照原始字段开始。若一项内容需要同时用编号和简称表示，整理表可保留两者，以便在不同视图中查找。为方便回查，可在任一表内记录有关的来源编号。这个示例仅展示字段之间的对应关系\NCRArticleCite{1}。

需要解释例外时，可在相应记录旁加入简短说明。正式文件中若没有提供可核对的材料，引用或图注应保留其原有范围。对照图时可用图号、面板标记及图例逐层核对。

\subsection*{方法示例}
先确认每个展室的记录来源，再将位置标记、观测时段和记录条件分别登记。将不同材料中的位置名称统一前，要检查其适用的时段。整理表可保存更新时间，供后续核对。

\subsection*{记录的可用性}
对可访问的资料，可列出访问路径和适用条件。作者声明、致谢及出版信息宜按各自标题排列。若某些附件仅见于原文引用，应按原载说明标识其范围\NCRArticleCite{2}。

\subsection*{记录的后续修订}
当位置发生变化时，需保留调整前的地点标记与相应记录。对相似的记录编号应加以区分，正文引用仍用同一个稳定的定位标记。纸面图注可另列交叉编号，便于读者从示意图回到原始材料。

对记录的更改应与源材料中作出的更订对应。文件仅调整列次时，可继续沿用原有编号。如果改用了新的观察位置，则应在整理表中并列写明原位置与新位置。查阅核查表的人可以据此知道记录适用的范围。

\subsection*{附录的连续阅读}
正文引用延伸至附表时，可在原条目旁标识表号与行号。若引用的图和完整图注因篇幅需要暂时分开，续页应提供可回查的标记。相邻文字可在图注之后接排，长表可按原有字段续页。

同一附录可能同时包含示例图、附表和表下注释。各区域的标题应与其实际内容对应，使后续编辑能定位需调整的部分。附录未提供具体条目时，可以如实注明材料范围。

\subsection*{核查记录}
检查记录时，可按表中顺序回看每项材料。把需要原始记录才能确认的问题单独标记，待补齐材料后与已核实的条目一并复核。不同版本重复描述的内容可各自注明出处。

随稿提供的附录可续排在相关文字之后。图注、表注、栏目说明和已填内容应分别对应原文。如果图的英文标签保留在原图内，译文可将中文释义放在相应图号下。

\section*{参考文献（示例占位）}
{\small\noindent\hypertarget{ncv-article:ref:1}{}1. Example, C. (20XX). Recording observational conditions. \emph{Example Journal}.\par
\noindent\hypertarget{ncv-article:ref:2}{}2. Example, D. (20XX). Reviewing linked entries. \emph{Example Journal}.\par}

\section*{资料与声明}
原文中的数据、代码及材料获取信息应分别照录。内部表格可随正文排版，能够独立取得的资源可保留链接及适用条件。各字段依原文所属的刊期或在线版本核对。

署名顺序、贡献分工、通信信息及许可条件可分别列出。示例作者的单位标记仅用来检验链接和字体，正式稿的单位名称与标记范围应逐项对应源文。
\end{multicols}

\medskip\hrule height .35bp\vspace{4bp}
\noindent{\sffamily\fontsize{8bp}{10.5bp}\selectfont
\hypertarget{ncv-article:affiliations}{}\textbf{作者单位（中性占位）}\enspace
\textsuperscript{1}第一示例单位。\enspace \textsuperscript{2}第二示例单位。\enspace
\textsuperscript{3}第三示例单位。\par
\noindent 共同贡献及通信信息按所给原稿填写。\par}

\section*{原文信息}
\noindent 作者、英文题名、期刊和 DOI\enspace 按所给原文填写。\par
\noindent 在线日期、正式卷期及页码\enspace 对应所给版本。此占位稿的日期、卷期和页码均未填。\par
\medskip\noindent 仅供学术交流，出版物版权归原作者及出版商所有。\par
\clearpage
\endgroup
